% GENERATED FILE. Edit canonical lessons, then run npm run generate:books.
% canonical-source-hash: fnv1a64:455e203b21fd236f
% canonical-lessons: BN-C506-tengk, BN-C506-phuto, BN-C506-upoche-pora, BN-C506-secha, BN-C506-chuiye-pora

\chapter{Leaks and Holes}
\label{ch:bn-leaks-and-holes}
\hlchaptermodality{506}

\begin{chapteropening}
\textbf{By the end of this chapter:} \emph{I can talk about water tanks, holes and water that seeps in.}

Complete the last lesson of chapter 506: I can talk about water tanks, holes and water that seeps in.
\end{chapteropening}

\section[\texorpdfstring{\bn{ট্যাংক} (ṭængk)}{ট্যাংক (ṭængk)}]{\bn{ট্যাংক} (ṭængk) — a water tank}

\label{lesson:BN-C506-tengk}



\begin{quote}
Before the new one: say the Bengali for to flicker, then the Bengali for to go out (of a light).
\end{quote}

\subsection*{You'll want to know: \bn{ট্যাংক}}
\textbf{\bn{ট্যাংক}} — \emph{ṭængk} — \textquotedblleft{}a water tank\textquotedblright{}.

\begin{quote}
\bn{ছাদের} \bn{ট্যাংকে} \bn{জল} \bn{নেই।}
\end{quote}

\begin{quote}
\emph{Chhāder ṭængke jôl nei.}
\end{quote}

\begin{quote}
There's no water in the tank on the roof.
\end{quote}

\begin{cousinweb}
\textbf{\bn{ট্যাংক}} (\emph{ṭængk}) is borrowed from English \emph{tank}.
\end{cousinweb}

\begin{grammarlens}[title={What you've built}]
One more word for describing what you lived and felt.
\end{grammarlens}

\subsection*{Your turn}
\begin{itemize}
\raggedright
  \item \emph{Say it:} \emph{ṭængk}
  \item \emph{Say it:} \emph{ṭængk}, once more
  \item \emph{Say it:} \emph{Chhāder ṭængke jôl nei.} — There's no water in the tank on the roof.
  \item \emph{Recall:} say the Bengali for to flicker, then the Bengali for to go out (of a light), then say \emph{ṭængk} again
\end{itemize}

\begin{tcolorbox}[breakable,colback=teal!4,colframe=teal!35!black,title={Before you move on}]
What does \bn{ট্যাংক} mean? (\textquotedblleft{}A water tank\textquotedblright{}.) Say it once more.
\end{tcolorbox}
\section[\texorpdfstring{\bn{ফুটো} (phuṭo)}{ফুটো (phuṭo)}]{\bn{ফুটো} (phuṭo) — a hole right through; full of holes}

\label{lesson:BN-C506-phuto}



\begin{quote}
Before the new one: say the Bengali for to go out (of a light), then the Bengali for a water tank.
\end{quote}

\subsection*{You'll want to know: \bn{ফুটো}}
\textbf{\bn{ফুটো}} — \emph{phuṭo} — \textquotedblleft{}a hole right through; full of holes\textquotedblright{}.

\begin{quote}
\bn{বালতিতে} \bn{একটা} \bn{ফুটো} \bn{আছে।}
\end{quote}

\begin{quote}
\emph{Bāltite ekṭā phuṭo āchhe.}
\end{quote}

\begin{quote}
There's a hole in the bucket.
\end{quote}

\begin{cousinweb}
\textbf{\bn{ফুটো}} (\emph{phuṭo}) belongs with a reconstructed \emph{sphuṭyati}, \textquotedblleft{}bursts open\textquotedblright{} (CDIAL 13845): a place where something has burst through.
\end{cousinweb}

\begin{grammarlens}[title={What you've built}]
One more word for describing what you lived and felt.
\end{grammarlens}

\subsection*{Your turn}
\begin{itemize}
\raggedright
  \item \emph{Say it:} \emph{phuṭo}
  \item \emph{Say it:} \emph{phuṭo}, once more
  \item \emph{Say it:} \emph{Bāltite ekṭā phuṭo āchhe.} — There's a hole in the bucket.
  \item \emph{Recall:} say the Bengali for to go out (of a light), then the Bengali for a water tank, then say \emph{phuṭo} again
\end{itemize}

\begin{tcolorbox}[breakable,colback=teal!4,colframe=teal!35!black,title={Before you move on}]
What does \bn{ফুটো} mean? (\textquotedblleft{}A hole right through; full of holes\textquotedblright{}.) Say it once more.
\end{tcolorbox}
\section[\texorpdfstring{\bn{উপচে} \bn{পড়া} (upôche pôṛā)}{উপচে পড়া (upôche pôṛā)}]{\bn{উপচে} \bn{পড়া} (upôche pôṛā) — to overflow; to brim over}

\label{lesson:BN-C506-upoche-pora}



\begin{quote}
Before the new one: say the Bengali for a water tank, then the Bengali for a hole right through.
\end{quote}

\subsection*{You'll want to know: \bn{উপচে} \bn{পড়া}}
\textbf{\bn{উপচে} \bn{পড়া}} — \emph{upôche pôṛā} — \textquotedblleft{}to overflow; to brim over\textquotedblright{}.

\begin{quote}
\bn{কলটা} \bn{বন্ধ} \bn{না} \bn{করলে} \bn{বালতি} \bn{থেকে} \bn{জল} \bn{উপচে} \bn{পড়বে।}
\end{quote}

\begin{quote}
\emph{Kôlṭā bôndhô nā kôrle bālti theke jôl upôche pôṛbe.}
\end{quote}

\begin{quote}
If you don't turn off the tap, the water will overflow the bucket.
\end{quote}

Past (he/she …): \bn{সেই} \bn{দুধ} \bn{উপচে} \bn{পড়ল}, \emph{sei dudh upôche pôṛlô}.

Said of milk boiling over too, as in the past form.

\begin{cousinweb}
\textbf{\bn{উপচে} \bn{পড়া}} (\emph{upôche pôṛā}) joins \bn{উপচে}, the -\bn{এ} form of a verb that CDIAL lists as Bengali \emph{upacāna}, \textquotedblleft{}to boil over, to overflow\textquotedblright{}, under a reconstructed \emph{utpacyatē}, \textquotedblleft{}boils over\textquotedblright{} (CDIAL 1808), to \bn{পড়া}, \textquotedblleft{}to fall\textquotedblright{}: the water falls over the rim as it spills.
\end{cousinweb}

\begin{grammarlens}[title={What you've built}]
One more word for describing what you lived and felt.
\end{grammarlens}

\subsection*{Your turn}
\begin{itemize}
\raggedright
  \item \emph{Say it:} \emph{upôche pôṛā}
  \item \emph{Say it:} \emph{upôche pôṛā}, once more
  \item \emph{Say it:} \emph{Kôlṭā bôndhô nā kôrle bālti theke jôl upôche pôṛbe.} — If you don't turn off the tap, the water will overflow the bucket.
  \item \emph{Recall:} say the Bengali for a water tank, then the Bengali for a hole right through, then say \emph{upôche pôṛā} again
\end{itemize}

\begin{tcolorbox}[breakable,colback=teal!4,colframe=teal!35!black,title={Before you move on}]
What does \bn{উপচে} \bn{পড়া} mean? (\textquotedblleft{}To overflow; to brim over\textquotedblright{}.) Say it once more.
\end{tcolorbox}
\section[\texorpdfstring{\bn{সেঁচা} (sẽchā)}{সেঁচা (sẽchā)}]{\bn{সেঁচা} (sẽchā) — to bail out water}

\label{lesson:BN-C506-secha}



\begin{quote}
Before the new one: say the Bengali for a hole right through, then the Bengali for to overflow.
\end{quote}

\subsection*{You'll want to know: \bn{সেঁচা}}
\textbf{\bn{সেঁচা}} — \emph{sẽchā} — \textquotedblleft{}to bail out water\textquotedblright{}.

\begin{quote}
\bn{বৃষ্টির} \bn{পরে} \bn{আমরা} \bn{উঠোনের} \bn{জল} \bn{সেঁচলাম।}
\end{quote}

\begin{quote}
\emph{Brishṭir pôre āmrā uṭhoner jôl sẽchlām.}
\end{quote}

\begin{quote}
After the rain we bailed the water out of the courtyard.
\end{quote}

Past (he/she …): \bn{তিনি} \bn{নৌকার} \bn{জল} \bn{সেঁচলেন}, \emph{tini noukār jôl sẽchlen}.

\begin{cousinweb}
\textbf{\bn{সেঁচা}} (\emph{sẽchā}) is a tadbhava: CDIAL lists Bengali \emph{sẽcā}, \textquotedblleft{}to irrigate, to bale out\textquotedblright{}, under Sanskrit \emph{siñcati}, \textquotedblleft{}pours out\textquotedblright{} (CDIAL 13394).
\end{cousinweb}

\begin{grammarlens}[title={What you've built}]
One more word for describing what you lived and felt.
\end{grammarlens}

\subsection*{Your turn}
\begin{itemize}
\raggedright
  \item \emph{Say it:} \emph{sẽchā}
  \item \emph{Say it:} \emph{sẽchā}, once more
  \item \emph{Say it:} \emph{Brishṭir pôre āmrā uṭhoner jôl sẽchlām.} — After the rain we bailed the water out of the courtyard.
  \item \emph{Recall:} say the Bengali for a hole right through, then the Bengali for to overflow, then say \emph{sẽchā} again
\end{itemize}

\begin{tcolorbox}[breakable,colback=teal!4,colframe=teal!35!black,title={Before you move on}]
What does \bn{সেঁচা} mean? (\textquotedblleft{}To bail out water\textquotedblright{}.) Say it once more.
\end{tcolorbox}
\section[\texorpdfstring{\bn{চুঁইয়ে} \bn{পড়া} (chũiye pôṛā)}{চুঁইয়ে পড়া (chũiye pôṛā)}]{\bn{চুঁইয়ে} \bn{পড়া} (chũiye pôṛā) — to seep; to leak through}

\label{lesson:BN-C506-chuiye-pora}



\begin{quote}
Before the new one: say the Bengali for to overflow, then the Bengali for to bail out water.
\end{quote}

\subsection*{You'll want to know: \bn{চুঁইয়ে} \bn{পড়া}}
\textbf{\bn{চুঁইয়ে} \bn{পড়া}} — \emph{chũiye pôṛā} — \textquotedblleft{}to seep; to leak through\textquotedblright{}.

\begin{quote}
\bn{বৃষ্টিতে} \bn{ছাদ} \bn{থেকে} \bn{জল} \bn{চুঁইয়ে} \bn{পড়ল।}
\end{quote}

\begin{quote}
\emph{Brishṭite chhād theke jôl chũiye pôṛlô.}
\end{quote}

\begin{quote}
In the rain, water seeped through the roof.
\end{quote}

Past (he/she …): \bn{সেই} \bn{জল} \bn{চুঁইয়ে} \bn{পড়ল}, \emph{sei jôl chũiye pôṛlô}.

\begin{cousinweb}
\textbf{\bn{চুঁইয়ে} \bn{পড়া}} (\emph{chũiye pôṛā}) joins \bn{চুঁইয়ে}, the -\bn{এ} form of a verb for oozing that CDIAL lists as Bengali \emph{cuyā}, \textquotedblleft{}to leak\textquotedblright{}, under Sanskrit \emph{cyotati}, \textquotedblleft{}trickles, oozes\textquotedblright{} (CDIAL 4948), to \bn{পড়া}, \textquotedblleft{}to fall\textquotedblright{}.
\end{cousinweb}

\begin{grammarlens}[title={What you've built}]
One more word for describing what you lived and felt.
\end{grammarlens}

\subsection*{Your turn}
\begin{itemize}
\raggedright
  \item \emph{Say it:} \emph{chũiye pôṛā}
  \item \emph{Say it:} \emph{chũiye pôṛā}, once more
  \item \emph{Say it:} \emph{Brishṭite chhād theke jôl chũiye pôṛlô.} — In the rain, water seeped through the roof.
  \item \emph{Recall:} say the Bengali for to overflow, then the Bengali for to bail out water, then say \emph{chũiye pôṛā} again
\end{itemize}

\begin{tcolorbox}[breakable,colback=teal!4,colframe=teal!35!black,title={Before you move on}]
What does \bn{চুঁইয়ে} \bn{পড়া} mean? (\textquotedblleft{}To seep; to leak through\textquotedblright{}.) Say it once more.
\end{tcolorbox}
